\scanpage{114}
\begin{center}- \mspage{122} -\end{center}
\margin{IX}
\[
(10.2.2)\qquad u_\ell^{A'_K}=u_\ell^{A_K}s\quad.
\]
Compte tenu de la propriété connue d'anticommutativité des accouplements $\varphi_\xi$ ([5] ou VIII 2.2.11), la formule de symétrie (10.2.2) peut se reformuler de la façon suivante, en termes d'une polarisation fixée $\xi$ de $A_K$, qui définit un homomorphisme $\widetilde\xi:A_K\longrightarrow A'_K$, d'où un homomorphisme correspondant ($t$ étant l'initiale de ``torique'')
\[
(10.2.3)\qquad\xi_t:\underline M\longrightarrow\underline M'\quad,
\]
permettant de former l'accouplement
\[
(10.2.4)\qquad\varphi_{\xi\ell}^t:\underline M_\ell\otimes\underline M_\ell\longrightarrow\mathbb Z_\ell,\quad
\varphi_{\xi,\ell}^t(x,y)=u_\ell^{A_K}(x,\xi_t(y))\quad:
\]
la forme précédente $\varphi_{\xi,\ell}^t$ est \uline{symétrique} :
\[
(10.2.5)\qquad\varphi_{\xi\ell}^t(x,y)=\varphi_{\xi\ell}^t(y,x)\quad.
\]
En fait, cet énoncé est valable pour tout élément $\xi$ du groupe de Néron-Séveri de $A_K$ (pas nécessairement une polarisation) ; mais ce n'est que lorsque $\xi$ est non dégénéré i.e. définit une \uline{isogénie} $A_K\longrightarrow A'_K$ donc aussi une isogénie (10.2.3), que la relation de symétrie (10.2.5) est équivalente à la relation de symétrie (10.2.2). D'ailleurs dans le cas où $\xi$ est une polarisation, on va préciser considérablement (10.2.5) dans 10.4 b) plus bas.

Signalons enfin une troisième façon d'interpréter la relation de symétrie. Pour ceci, considérons une biextension (= correspondance divisorielle) $\xi$ sur un couple $(A_K,\overline A_K)$ de schémas abéliens à réduction semi-stable sur $K$, qui définit un homomorphisme
\[
A_K\longrightarrow\overline A'_K\quad,
\]
\margin{440}
\margin{440}
\margin{113}
